\chapter{Prezentare aplicație}
\section{Sign in si sign out}
Aplicația cere fiecărui utilizator să își creeze un cont pe baza unui email și a unei parole pentru folosirea funcționalităților. Fiecare cont are nevoie de confirmarea email-ului, abia după confirmare utilizatorul se poate loga și folosi aplicația.

Ecranele de sign in și sign out au atașate validări împreună cu mesaje specifice pentru fiecare caz: format incorect pentru email și parolă, parolă incorectă, user inexistent. Aplicația oferă posibilitatea deconectării și ștergerii contului, având un mesaj de avertizare că datele se vor șterge, cu verificarea parolei înainte, pentru ca ștergerea să poată fi făcută doar de utilizator.

\section{Pagina principala}
Pagina principală (home) este concepută atât ca punct de intrare pentru vizitatori, cât și ca panou de control pentru utilizatorii autentificați. Pentru utilizatorii neautentificați prezintă o pagină cu detalii despre aplicație împreună cu butoane de logare și înregistrare. După logare, pagina prezintă un mesaj de bun venit și butoane care ghidează utilizatorul să creeze o nouă călătorie sau să le vizualizeze pe cele existente.

\section{Călătoriile mele}
„My Trips" (Călătoriile mele) este pagina centrală pentru gestionarea călătoriilor unui utilizator. Călătoriile sunt separate în trei categorii distincte, în funcție de statutul lor temporal. Secțiunea „Se întâmplă acum" (Happening Now) apare doar dacă o excursie este în desfășurare și afișează călătoriile active în prezent pentru un acces mai rapid la acestea. Secțiunea „Aventuri viitoare" prezintă călătoriile care urmează în viitor în ordine calendaristică pentru a ajuta utilizatorul să vizualizeze călătoriile care urmează și pentru a se pregăti pentru ele. Ultima secțiune, „Amintiri de călătorie" (Travel Memories), păstrează câteva călătorii terminate, permițând utilizatorului să vizualizeze excursiile trecute împreună cu detaliile acestora, având la dispoziție un buton de „see more" pentru a accesa o pagină cu toate călătoriile terminate.

Fiecare călătorie apare ca un card interactiv care afișează informațiile cheie, cum ar fi data și destinația. În cazul primelor două secțiuni, cardurile au atașate butoane către adăugarea de atracții, cât și un buton pentru generarea de itinerariu și vederea tuturor detaliilor. În ultima secțiune butonul de adăugare este eliminat și rămâne doar cel pentru vizualizarea detaliilor.

Un buton „Create new trip" (Creează o nouă călătorie) permite utilizatorilor să acceseze ușor crearea unei noi călătorii.

\section{Creare călătorie}
Pagina „Plan new trip" (Creează o nouă călătorie) conține un formular complet prin care utilizatorii pot crea o nouă călătorie. În formular utilizatorii introduc informații precum numele, o descriere în cazul în care aceștia doresc, destinația, datele de început și sfârșit. De asemenea formularul include selectoare pentru orele pe care utilizatorul ar vrea să le folosească pentru vizitare, folosite la generarea itinerariului. Validarea formularului asigură integritatea datelor înainte de terminare, ajutând utilizatorul să completeze corect câmpurile. Această pagină este primul pas în procesul de planificare a călătoriei și este baza pentru următorii pași.

\section{Detaliile călătoriei}
Pagina „Details" (Detalii) este cea mai cuprinzătoare și include cele mai multe funcții din aplicație, concepută pentru a oferi o imagine de ansamblu completă și o interfață de gestionare a unei călătorii. Pagina prezintă informația defalcată în două secțiuni pentru o înțelegere mai clară a acesteia.

Coloana din stânga prezintă toate detaliile călătoriei selectate. În partea de sus avem o listă colapsabilă cu toate locațiile adăugate și cele două butoane de generare. Opțional, utilizatorul poate adăuga adresa cazării în care va sta, care va apărea pe hartă pentru o mai bună generare a itinerariului. Odată ce itinerariul este generat, în secțiunea „Daily itinerary" se afișează o listă cu fiecare zi din excursia selectată și un orar cu locațiile asociate acelei zile. Fiecare locație are o casetă de selectare unde utilizatorul poate bifa locațiile vizitate.

Coloana din dreapta prezintă o hartă interactivă generată prin Mapbox în care putem vizualiza traseul geografic al călătoriei. În această hartă sunt afișate toate locațiile prin pinuri pe hartă. După generarea traseului, aceste locații sunt conectate prin rute de culori diferite în funcție de ziua din care fac parte, pentru o mai bună vizualizare spațială și temporală a fiecărei zile.

\subsection{Generare simpla a itinerariului}
Funcția „GenerateItinerary" (Generare itinerariu) este construită pe o euristică Greedy Nearest-Neighbor modificată, o abordare derivată din Problema Comis Voiajorului (TSP). În forma canonică, algoritmul Nearest-Neighbor începe de la un nod dat și selectează în mod repetat cel mai apropiat nod nevizitat până când toate nodurile au fost vizitate. În implementarea noastră, acest concept a fost adaptat pentru a funcționa în cadrul constrângerilor creării unui itinerariu de călătorie pe mai multe zile, unde ferestrele de timp, rezervările, disponibilitatea locației și preferințele utilizatorului trebuie respectate simultan.

\ding{117}{Ce am modificat}

Algoritmul standard Nearest-Neighbor funcționează exclusiv pe distanță, alege cel mai apropiat punct nevizitat geografic la fiecare pas. Versiunea noastră înlocuiește distanța geografică brută cu timpul real de călătorie, obținut din API-ul Mapbox Directions Matrix, care returnează duratele reale de mers pe jos și cu mașina între poziția curentă și toate destinațiile candidate. La fiecare pas algoritmul încearcă să găsească cea mai apropiată locație din punct de vedere temporal, încadrându-se în fereastra de timp disponibilă. Această schimbare este importantă deoarece două locații echidistante pe hartă pot avea timpi diferiți de călătorie în funcție de drumurile parcurse.

Deoarece în aplicația noastră apar restricții de timp, vom introduce un concept de constrângeri temporale având atât rezervări (rezervări externe fixe cu ore de început și sfârșit), cât și elemente vizitate (locații pe care utilizatorul le-a vizitat). În faza de curățare, elementele vizitate sunt relocate, astfel încât următoarele zile să fie reconstruite de la zero. Într-o anumită zi, la fiecare iterație, algoritmul analizează în perspectivă următorul eveniment fix. Înainte de a plasa o locație nouă în program, algoritmul verifică dacă utilizatorul ar avea suficient timp pentru a ajunge acolo, a petrece durata medie și a ajunge la următorul eveniment cu o marjă de eroare de 15 minute. Se încearcă până se găsește o locație care poate respecta aceste condiții, pentru a evita crearea unui conflict cu o rezervare preexistentă.

\ding{117}{Ce folosește algoritmul pentru a calcula}

Fiecare iterație procesează fiecare DayPlan cronologic. Pentru fiecare zi, bucla inițiază un cursor currentTime la ora configurată de începere a explorării și un cursor endTimeLimit la ora de sfârșit a explorării. Dacă cazarea a fost adăugată, aceasta este ancorată ca poziție de începere, astfel încât prima iterație în fiecare zi începe de la cazare. Bucla apelează apoi FindBestNextLocationAsync.

În interiorul acestei funcții, se verifică dacă există mai mult de 24 de locații candidate. Dacă da, acestea sunt prefiltrate folosind formula Haversine, care calculează distanța dintre două coordonate geografice ținând cont de curbura planetei, pentru a selecta doar cele 24 de locații cele mai apropiate. Aceasta este o optimizare necesară deoarece API-ul Mapbox Matrix are o limită de 25 de puncte per solicitare, astfel încât prefiltrul Haversine acționează ca o modalitate ușoară și rapidă de a elimina locațiile prea îndepărtate înainte de a efectua apelul API. Calculul Haversine convertește diferențele de latitudine și longitudine în radiani, aplică legea sferică a cosinusurilor prin termenii sin² și cos și returnează distanța în metri prin raza Pământului de 6.371.000 de metri.

Odată ce lista de candidați a fost redusă, funcția apelează API-ul Mapbox Matrix de două ori în paralel, o dată pentru durata de mers pe jos și o dată pentru durata de condus între poziția curentă și toți candidații. Pentru fiecare candidat algoritmul alege care este cel mai rapid mijloc de a ajunge în acea locație. Dacă API-ul eșuează sau nu returnează niciun rezultat pentru un anumit candidat, algoritmul revine la un timp de călătorie estimat calculat din distanța Haversine împărțind la o viteză medie de 10 km/h, aproximând un ritm de mers rapid. Această soluție asigură că algoritmul nu se oprește din cauza defecțiunilor rețelei.

Pentru fiecare candidat, algoritmul aplică două filtre înainte de a-l considera cea mai bună alegere. Primul filtru este IsLocationOpen, care verifică orele de deschidere și închidere ale locației în raport cu ora de sosire proiectată. Al doilea este FitsBeforeNextEvent, care verifică dacă vizita completă a locației, sosirea plus durata medie, se încheie în timp util pentru a ajunge la următoarea locație; în cazul rezervărilor această constrângere este mărită cu 15 minute. Candidații care trec de ambele filtre sunt luați în considerare, fiind ales cel care are cel mai scurt timp de călătorie, păstrând regulile algoritmului Greedy.

Algoritmul rezolvă și problema schimbărilor în itinerariu, atunci când utilizatorul nu reușește să viziteze totul dintr-o anumită zi sau înregistrează o locație înainte de program. În loc să permită blocarea programului unei întregi zile, detectează locațiile care au fost vizitate înainte și le mută în ziua curentă din baza de date. Istoricul locației este păstrat, dar ziua viitoare este eliberată astfel încât itinerariul să poată fi regenerat. După curățare, algoritmul construiește setul de locații încă disponibile pentru programare prin interogarea elementelor itinerariului și colectarea celor marcate „Vizitate". Acestea, împreună cu orice rezervare, sunt ulterior excluse din grupul de candidați, rămânând doar locațiile care pot fi reprogramate.

Această strategie de regenerare face ca algoritmul să fie idempotent și sigur pentru a putea fi rulat din nou. Nu există riscul ștergerii istoricului sau completării retroactive a unei zile deja trecute.

\subsection{Generare cu ai}
Funcția GenerateWithAI implementează un model fundamental diferit: nu este un algoritm de rutare în sine, ci mai degrabă o conductă în două etape care combină un LLM pentru descoperirea semantică a locurilor cu API-ul de Geocodare Mapbox pentru validarea coordonatelor din lumea reală, urmată de un apel recursiv către GenerateItinerary pentru a programa efectiv locațiile nou găsite.

\ding{117}{Solicitarea LLM bazate pe preferințe}

Prima etapă interoghează preferințele de etichete stocate ale utilizatorului din baza de date, preluând categoriile pe care le-a clasificat. Acestea sunt transmise într-o solicitare structurată trimisă modelului LLaMA 3.3 70B prin intermediul API-ului Groq. Modelul este solicitat cu o instrucțiune de sistem constrânsă: trebuie să răspundă doar într-o schemă JSON specifică care conține o matrice de obiecte, fiecare cu un nume și o etichetă. Partea de solicitare orientată către utilizator solicită până la 8 locuri de interes reale, vizitabile, în orașul și țara călătoriei sau în apropiere, prioritizate în funcție de etichetele principale ale utilizatorului și exclude în mod explicit hoteluri, restaurante și cartiere.

Răspunsul modelului este analizat, iar șirul JSON brut este curățat prin găsirea primelor caractere `{}` și extragerea datelor, măsură care asigură evitarea oricărui preambul sau text final pe care modelul l-ar putea produce în ciuda solicitării stricte.

\ding{117}{Validarea}

Deoarece LLM-ul poate halucina unele nume sau poate returna locații din orașul greșit, fiecare sugestie returnată de model este validată independent prin intermediul API-ului Mapbox Search Forward Geocoding. Pentru fiecare nume de loc sugerat, funcția construiește o interogare de geocodare care combină numele locului cu orașul călătoriei și adaugă o verificare de proximitate folosind coordonatele centrului geografic al orașului. Această verificare este necesară pentru a returna doar rezultate din interiorul orașului sau din apropierea orașului, mai degrabă decât un loc cu același nume din altă parte.

Fiecare rezultat de geocodare este supus unor filtre de validare. Numele este verificat în raport cu o listă de cuvinte cheie precum „hotel", „apartament", „parcare" și altele asemenea, pentru a nu intra în itinerariu. Adresa rezultatului este verificată pentru a confirma că menționează țara și orașul corect. Dacă numele orașului nu apare în adresă, se aplică o verificare de proximitate prin formula Haversine pentru a calcula distanța dintre coordonatele rezultatului și centrul orașului, iar rezultatele aflate pe o rază de 30 de km sunt acceptate chiar dacă numele lipsește din textul adresei. Aceasta gestionează cazurile în care o excursie este planificată într-o locație mică, unde locațiile ar putea fi situate adiacent. În cele din urmă, rezultatul este verificat atât în raport cu locațiile deja existente, cât și cu cele deja adăugate, pentru a preveni intrările duplicate. Pentru fiecare locație se creează o nouă entitate Locație care se adaugă în baza de date.

Odată ce toate locațiile au fost salvate, GenerateWithAI se termină cu un apel către GenerateItinerary, deoarece funcția AI este un mecanism de descoperire și populare cu locații, în timp ce algoritmul propriu-zis de generare se află exclusiv în GenerateItinerary. Cele două generări au în comun un singur motor de programare unificat, astfel încât îmbunătățirile uneia o avantajează și pe cealaltă.

\section{Adăugare obiective}
„Add Location" (Adăugare locații) permite utilizatorilor să adauge puncte de interes în itinerariul călătoriei alese. În stânga se află harta interactivă unde fiecare locație adăugată este evidențiată printr-un pin colorat specific etichetelor din care face parte. De asemenea se pot selecta manual locații direct din hartă și ulterior se pot adăuga în listă, sau se pot căuta orice locații folosind bara de căutare de deasupra hărții. Aceasta folosește cuvinte cheie introduse și oferă variante care se potrivesc. Odată ce o locație este selectată, pinul va apărea pe hartă, iar în dreapta secțiunea „Add new" se va completa cu datele acelei locații. Pentru fiecare locație sunt prezentate numele și locația, o listă de etichete care îi sunt asociate. De asemenea avem câmpuri suplimentare cu durata aproximativă de vizitare și orele în care această locație este deschisă, acestea autocompletându-se dar cu posibilitatea de editare. Apăsând pe butonul „Add to itinerary" (Adaugă în itinerariu), pinul se colorează corespunzător și locația rămâne salvată în listă.

Secțiunea de sugestii este o funcție de recomandări concepută pentru a ajuta utilizatorii să descopere atracții și puncte de interes din orașul pe care îl vizitează. Meniul este organizat pe categorii, cu un meniu derulant care oferă diverse opțiuni, cum ar fi Muzee, Parcuri, Restaurante, Atracții și multe altele. Utilizatorii pot alege variante aleatoare pentru a obține o colecție diversă de locuri sau pot alege o categorie specifică. Când o locație este accesată, toate specificațiile vor apărea în pagina de adăugare.

Ultima secțiune este lista de locații pe care utilizatorul le are adăugate în călătoria curentă. Aceasta este o vizualizare rapidă a tuturor locațiilor, fiecare apărând într-un card interactiv care afișează toate detaliile. Dacă utilizatorul apasă pe un card, harta se va centraliza pe acea locație și detaliile vor apărea în pagina de adăugare. Fiecare card are un buton de ștergere și unul de „Add Event" (Adaugă eveniment), care dacă este apăsat va deschide un card unde se pot adăuga detaliile unui eveniment care va avea loc la acea locație: numele, ora începerii și durata.

\section{Pagina calendar}
Pagina „Calendar" (Calendar) prezintă un cuprins al tuturor călătoriilor unui utilizator într-un format de calendar, oferind o imagine de ansamblu rapidă asupra programului de călătorie pe tot parcursul anului. Fiecare călătorie este evidențiată, facilitând identificarea momentului când are loc o călătorie și planificarea în funcție de acestea. Utilizatorii pot naviga între vizualizare pe luni sau ani, iar apăsând pe o anumită călătorie sunt redirecționați rapid la pagina cu detalii corespunzătoare. Această vizualizare ajută utilizatorii să își organizeze mai bine excursiile pe parcursul anului și oferă o imagine de ansamblu a acestora.

\section{Harta interactivă}
„ScratchMap" (Harta Răzuibilă) este o funcție de gamificare care afișează o hartă interactivă a lumii ce evidențiază toate țările vizitate anterior și cele care vor urma. Pe măsură ce utilizatorii creează și finalizează călătorii, țările sunt evidențiate automat în culori distincte pe hartă, verde pentru cele terminate și galben pentru cele care vor urma, creând o imagine de ansamblu vizuală a explorării globului. Această funcție oferă o vizualizare captivantă a istoricului călătoriilor, aliniindu-se cu obiectivul aplicației de a face planificarea călătoriilor mai plăcută.

\section{Pagina de preferințe}
În pagina „My preferences" (Preferințele mele) utilizatorul își poate configura interesele și prioritățile, care servesc ca bază pentru algoritmul de recomandare bazat pe inteligența artificială. Pagina prezintă o colecție cuprinzătoare de categorii și etichete de interes organizate pe trei niveluri de prioritate. Interacțiunea cu acestea se face printr-o interfață intuitivă care permite glisarea și plasarea etichetelor în unul din cele trei niveluri de prioritate. Acest sistem este util în generarea itinerariilor, algoritmul ponderând recomandările în funcție de aceste priorități, asigurându-se că planurile de călătorie se aliniază cu interesele fiecărui utilizator.
